\chapter{Reciprocal--Conjugation Orbit Collapse in the Projective Coordinate}
\label{chap:reciprocal-conjugation-orbit-collapse}

\status{Exact RH-equivalent reformulation and orbit-separating geometry. No implication from zerohood to orbit collapse is assumed or proved here; RH remains OPEN.}

\section{Projective coordinate and transformed xi function}
Set
\[
u=\Omega(s)=\frac{s}{1-s},\qquad s=\frac{u}{1+u},
\]
and define
\[
X(u)=\xi\!\left(\frac{u}{1+u}\right).
\]
For a non-trivial xi zero $\rho$, write $u_\rho=\Omega(\rho)$. Since non-trivial zeros lie in the open critical strip, $u_\rho\ne0,\infty$.

\section{Two exact zero-set actions}
The holomorphic functional reflection
\[
\mathcal I_s(s)=1-s
\]
becomes
\[
\boxed{\mathcal I_u(u)=\frac1u.}
\]
Because $\xi(s)=\xi(1-s)$,
\[
X(u)=0\Longrightarrow X(1/u)=0.
\]

Complex conjugation gives
\[
\boxed{\mathcal C_u(u)=\overline u,}
\]
and, because $\xi(\bar s)=\overline{\xi(s)}$,
\[
X(u)=0\Longrightarrow X(\bar u)=0.
\]
These two involutions commute:
\[
\mathcal I\mathcal C(u)=\frac1{\bar u}
=\mathcal C\mathcal I(u).
\]
Thus every transformed zero belongs to an orbit contained in
\[
\boxed{\left\{u,\bar u,\frac1u,\frac1{\bar u}\right\}.}
\]

\section{Critical line as equality of the two actions}
\begin{theorem}[Reciprocal--conjugation coincidence]
For every $u\in\mathbb C^\times$,
\[
\boxed{
\mathcal I_u(u)=\mathcal C_u(u)
\iff
|u|=1.
}
\]
Under $u=\Omega(s)$ this is equivalent to $\Re s=1/2$.
\end{theorem}

\begin{proof}
The equality is
\[
\frac1u=\bar u.
\]
Multiplication by $u$ gives $u\bar u=1$, hence $|u|=1$, and the converse is immediate. The projective identity
\[
|\Omega(s)|=1\iff |s|=|1-s|\iff\Re s=\frac12
\]
completes the crosswalk.
\end{proof}

Consequently,
\[
\boxed{
\mathrm{RH}
\iff
\forall u\in Z_X:\ \mathcal I_u(u)=\mathcal C_u(u).
}
\]
This is an exact equivalence, not an independent proof of the implication.

\section{Orbit cardinality}
For a non-trivial zero, the generic off-axis orbit has four distinct points
\[
\left\{u,\bar u,\frac1u,\frac1{\bar u}\right\}.
\]
On the unit circle,
\[
\frac1u=\bar u,
\]
so the reciprocal and conjugation pairings coincide and the orbit reduces to the conjugate pair
\[
\{u,\bar u\}.
\]
The exceptional fixed values of either involution are not used as a zero-location premise. In particular, pointwise fixedness under $u\mapsto1/u$ would force $u=\pm1$ and is not the RH condition.

\section{Orbit-separating defect}
Define
\[
\boxed{
\Delta_{\mathrm{RC}}(u)
:=\left|\frac1u-\bar u\right|^2.
}
\]
Writing $R=|u|$ gives the exact identities
\[
\Delta_{\mathrm{RC}}(u)
=\frac{(1-R^2)^2}{R^2}
=\left(R-\frac1R\right)^2.
\]
Hence
\[
\boxed{\Delta_{\mathrm{RC}}(u)=0\iff R=1.}
\]

For the compactified projective radius
\[
q=\frac{R}{1+R},
\]
one has
\[
R=\frac{q}{1-q}
\]
and therefore
\[
\boxed{
\Delta_{\mathrm{RC}}(q)
=\frac{(2q-1)^2}{q^2(1-q)^2}.
}
\]
Thus
\[
\Delta_{\mathrm{RC}}=0
\iff q=\frac12
\iff\Re s=\frac12.
\]

For the logarithmic radius
\[
B=\log R,
\]
\[
\boxed{\Delta_{\mathrm{RC}}=4\sinh^2B.}
\]
For the centred compact coordinate
\[
\eta=\frac{R-1}{R+1}=2q-1,
\]
\[
\boxed{
\Delta_{\mathrm{RC}}
=\frac{16\eta^2}{(1-\eta^2)^2}.
}
\]
All four forms vanish on exactly the same self-dual layer.

\section{Why orbit averaging cannot close the problem}
Reciprocal inversion sends
\[
q\mapsto1-q,
\qquad B\mapsto-B.
\]
Every off-axis two-cycle therefore satisfies
\[
\frac{q+(1-q)}2=\frac12,
\qquad B+(-B)=0.
\]
Consequently, an orbit barycentre or an odd orbit sum cannot distinguish a genuine fixed/self-dual layer from a nontrivial two-cycle. The squared defect $B^2$ and $\Delta_{\mathrm{RC}}$ retain the pointwise separation.

\section{Minimal remaining implication}
In these coordinates the remaining RH obligation can be written without auxiliary machinery as
\[
\boxed{
X(u)=0\Longrightarrow\Delta_{\mathrm{RC}}(u)=0.
}
\]
Equivalently,
\[
X(u)=0\Longrightarrow |u|=1,
\]
or
\[
X(u)=0\Longrightarrow q(u)=\frac12.
\]
This implication is RH-equivalent. The exact orbit algebra above compresses the target but does not supply the required non-circular incoming theorem.

\section{Separation from negative inversion}
The map used in this chapter is the functional-equation action
\[
\mathcal I_u(u)=1/u.
\]
It must not be conflated with the Li/Euler negative inversion
\[
\mathcal N_u(u)=-1/u.
\]
Chapters~37--43 prove exact geometry for $\mathcal N$ and then exclude it as a zero-to-zero spectral mechanism. Those no-go results do not invalidate the functional symmetry $u\mapsto1/u$ used here.

\gap{The orbit-collapse equivalence and defect formulas are exact. The unproved statement is precisely $X(u)=0\Rightarrow\Delta_{\mathrm{RC}}(u)=0$. No QED for RH is claimed.}
